\documentclass[10pt,a4paper]{amsart}
\usepackage[margin=19mm]{geometry}
\usepackage{amsmath,amssymb,mathtools}
\usepackage[scr=rsfs,cal=euler]{mathalfa}
\usepackage{xfrac}
\usepackage[unicode,hidelinks,bookmarks=false]{hyperref}
\newcommand{\ie}{i.e.\ }
\newcommand{\cf}{cf.\ }
\newcommand{\resp}{respectively\ }
\newcommand{\Subsection}{\subsection}
\providecommand{\nobreakdash}{\nobreak\hbox{-}}
\setlength{\emergencystretch}{2em}
\multlinegap0pt
\usepackage{bm}
\usepackage{bbm}
\usepackage{dsfont}
\usepackage{xspace}
\usepackage{cancel}
\usepackage{mathtools}
\usepackage{amsthm}
\makeatletter
\@ifundefined{theorem}{\newtheorem{theorem}{Theorem}}{}
\@ifundefined{proposition}{\newtheorem{proposition}{Proposition}}{}
\makeatother
\makeatletter
\newcommand{\Fabrice}[1]{#1}
\expandafter\ifx\csname del\endcsname\relax
\newenvironment{del}
  {}
  {}
\fi
\makeatother
\title[An Asymptotic-Preserving Micro--Macro Scheme for Plasma Simu]{An Asymptotic-Preserving Micro--Macro Scheme for Plasma Simulations in Quasi-Neutral and Low-Mach-Number Regimes with Kinetic Upgrades}
\author{Zeyu Liu, Fabrice Deluzet, Chang Yang}
\date{Source: arXiv:2608.07197v1 (CC BY 4.0)}
\begin{document}
\maketitle

\noindent\emph{Source and licence.} An Asymptotic-Preserving Micro--Macro Scheme for Plasma Simulations in Quasi-Neutral and Low-Mach-Number Regimes with Kinetic Upgrades. Version arXiv:2608.07197v1 (CC BY 4.0), distributed under the Creative Commons Attribution 4.0 International licence, as recorded in the arXiv licence metadata for this exact version and shown on the abstract page https://arxiv.org/abs/2608.07197v1. Full rights notice: licenses/arxiv-2608.07197-CC-BY-4.0.txt.\par\smallskip

\noindent\textit{Editorial scope.} Counted unit: the Proposition whose source label is \texttt{prop2.2} in the article (source0.tex lines 444--460, source label \texttt{prop2.2}), delivered together with the article's own complete original proof of it (source0.tex lines 462--493, one \texttt{proof} environment of 32 lines). No mathematics is written, repaired or re-derived: statement and proof are the article's own text line for line. Required context retained uncounted: MMd (lines 399--411); MMm (lines 401--405). Every definition, display and label of those lines is retained unchanged. Omitted from this package and not redistributed: 1--398, 406--443, 461--461, 494--2217. Nothing inside a retained excerpt is omitted or altered: every retained body is delivered exactly as the article's lines are written, so no omission receipt of any kind accompanies this package. Citations: the retained excerpts carry no citation command at all, so no bibliography entry of the article is transcribed and no reference is redistributed. Reference adaptation: none was needed and none was made; every reference inside the delivered unit resolves inside it, so no unresolved reference or citation is produced. Printed slips kept literal: no printed word, symbol, hypothesis or number of the counted statement or of its proof was altered. Layout and class: the article's own class is \texttt{article} with packages including epsfig, graphics, latexsym, verbatim, amsmath, amssymb, comment, units, amsfonts, bm, graphicx, subcaption, placeins, booktabs; the wrapper is the lane's pinned \texttt{amsart} editorial wrapper at 10pt on A4, which restates the theorem-like environments the retained text needs and adds the article's own macro definitions that the retained text needs (\texttt{\textbackslash Fabrice}), with extra wrapper packages (bm, bbm, dsfont, xspace, cancel, mathtools). The delivered numbering is the unit's own. Assets: no retained span contains a figure environment, a graphics-inclusion command, a PDF-page inclusion, a table or a TikZ picture; no image file is copied into this package, the package declares no asset dependency and no image file is redistributed with it. Licence: version arXiv:2608.07197v1 of this article is distributed under the Creative Commons Attribution 4.0 International licence, as recorded in the arXiv licence metadata for this exact version and shown on the abstract page https://arxiv.org/abs/2608.07197v1. A full rights notice is delivered at licenses/arxiv-2608.07197-CC-BY-4.0.txt. Characters: every retained excerpt is pure ASCII, so the delivered engine, XeLaTeX with its default Latin Modern text font, reports no missing character.
\begin{align}
    &\partial_{t}{n_e}+\nabla_{x}\cdot(n_e u_e)=0,\label{MMd}\\
    \begin{split}
    &M^{2}(\partial_{t}(n_e u_e)+\nabla_{x}\cdot(n_e u_e\otimes u_e))+T_{e,0}\nabla_{x}n_e-
    \eta n_e\nabla_{x}\phi\\
    &\hspace{5.5cm}+\Omega_{e}M^{2}n_e u_e\times B = -\nabla_{x}\cdot\langle v\otimes vg_{e}\rangle,
    \end{split}\label{MMm}\\
    \begin{split}
    &M\partial_t{g_e}+(\mathbb{I}-\Pi_{\mathcal{M}})(v\cdot\nabla_xg_e-\Omega_eMv\times B\cdot\nabla_vg_e)-\eta E\cdot\nabla_vg_e \\
    &\hspace{2cm}=-\frac{M}{\kappa}\nu_{ee}g_e-(\mathbb{I}-\Pi_{\mathcal{M}})(v\cdot\nabla_x\mathcal{M}_{e}-\Omega_eMv\times B\cdot\nabla_v\mathcal{M}_{e}),
    \end{split}\label{MMk}\\
    &-\lambda^{2}\eta\Delta_{x}\phi=n_{i}-n_e,\label{MMp}
\end{align}

\begin{proposition}[Reformulated Poisson equation and quasi-neutral model] \label{prop2.2}
Within the micro--macro framework, assuming that the initial conditions satisfy both Poisson equation and its first-order time derivative, namely:
\begin{equation} \label{eq:init_cond}
    -\lambda^2\eta\Delta_x\phi\vert_{t=0} = n_i - n_e\vert_{t=0} \quad \text{and} \quad -\lambda^2\eta\Delta_x\partial_t\phi\vert_{t=0} = \nabla_x\cdot(n_e u_e)\vert_{t=0},
\end{equation}
then Poisson equation is equivalent for all $t > 0$ to the following reformulated equation:
\begin{align} \label{eq:poisson_reformulated}
    -\nabla_x\cdot\Big((\Fabrice{\eta}n_e+M^2\lambda^2\eta\partial_t^2)\nabla_x\phi\Big) = &- M^2\nabla_x^2:(n_eu_e\otimes u_e)-T_{e,0}\Fabrice{\Delta_x}n_e \\ &\Fabrice{ - \Omega_e M^2 \nabla_x \cdot (n_e u_e \times B )} - \nabla_x^2:\left< v\otimes vg_e \right>. \notag
\end{align}
where $\nabla_x^2 : T = \nabla_x \cdot (\nabla_x \cdot T)$ denotes the double divergence of a second-order tensor $T$.

In the limit $\lambda\to0$, the reformulated Poisson equation yields the following elliptic equation for the electric potential, subject to suitable boundary conditions:
\begin{align*}
    -\nabla_x\cdot(\Fabrice{\eta}n_e\nabla_x\phi) = &- M^2\nabla_x^2:(n_eu_e\otimes u_e)-T_{e,0}\Fabrice{\Delta_x}n_e \\
   &\quad \Fabrice{ - \Omega_e M^2 \nabla_x \cdot (n_e u_e \times B )} - \nabla_x^2:\left< v\otimes vg_e \right>.
\end{align*}
\end{proposition}

\begin{proof} \label{prop2.2pf}
The proof consists in time differentiating the continuity equation $\partial_t \rho + \nabla_x \cdot J=0$ to establish a wave-like equation for the charge density, which is then combined with the second-order time derivative of Poisson equation. Under the assumption of stationary ions, the continuity equation reduces to the electron density conservation law.

Differentiating equation \eqref{MMd} with respect to time and taking the divergence of the momentum equation \eqref{MMm} yields:
\begin{align*}
        &\partial_t^2n_e + \partial_t\nabla_x\cdot(n_eu_e) = 0, \\
        &M^2\partial_t\nabla_x\cdot(n_eu_e) + M^2\nabla_x^2:(n_eu_e\otimes u_e)+T_{e,0}\Fabrice{\Delta_x}n_e  = \\
        & \hspace{3.5cm} \Fabrice{\eta} \nabla_x\cdot(n_e\nabla_x\phi)\Fabrice{ - \Omega_e M^2 \nabla_x \cdot (n_e u_e \times B )}-\nabla_x^2:\left< v\otimes vg_e \right>.
\end{align*}
Eliminating the cross-derivative term $\partial_t\nabla_x\cdot(n_eu_e)$ between these two relations yields:
\begin{equation} \label{eq:wave_ne}
    \begin{split}
        M^2\partial_t^2n_e + \Fabrice{\eta}\nabla_x\cdot(n_e\nabla_x\phi) &= M^2\nabla_x^2:(n_eu_e\otimes u_e)+T_{e,0}\Fabrice{\Delta_x}n_e \\
        &\quad\Fabrice{ + \Omega_e M^2 \nabla_x \cdot (n_e u_e \times B )} + \nabla_x^2:\left< v\otimes vg_e \right>.
    \end{split}
\end{equation}
The reformulated Poisson equation is obtained by substituting $M^2\partial_t^2n_e$ into the double time derivative of Poisson equation: $-M^2 \lambda^2 \partial_t^2 \Delta_x \phi = -M^2 \partial_t^2 n_e$ providing equation \eqref{eq:poisson_reformulated}.

Now, let us define the deviation from Poisson equation as $\mathcal{E}(t,x) = \lambda^2 \eta \Delta_x \phi + n_i-n_e$, which vanishes if Poisson equation holds true. Since the ions are at rest ($\partial_t^2 n_i = 0$), the second-order time derivative of $\mathcal{E}$ reads $M^2 \partial_t^2 \mathcal{E} = -M^2 \partial_t^2 n_e + M^2 \lambda^2 \eta \partial_t^2 \Delta_x \phi$. Substituting $M^2 \partial_t^2 n_e$ from \eqref{eq:wave_ne} into this relation yields:
\begin{equation}\label{eq:Poisson:deviation}
    \begin{multlined}[0.9\textwidth]
    M^2\partial_t^2 \mathcal{E} = \nabla_x\cdot\Big((\Fabrice{\eta}n_e+M^2\lambda^2\eta\partial_t^2)\nabla_x\phi\Big) \\ - M^2\nabla_x^2:(n_eu_e\otimes u_e)-T_{e,0}\Fabrice{\Delta_x}n_e \Fabrice{ - \Omega_e M^2 \nabla_x \cdot (n_e u_e \times B )} - \nabla_x^2:\left< v\otimes vg_e \right>.
    \end{multlined}
\end{equation}
Enforcing the reformulated Poisson equation \eqref{eq:poisson_reformulated} is mathematically equivalent to canceling the source term in \eqref{eq:Poisson:deviation}. Since this equation dictates the dynamics of the deviation, we obtain the homogeneous relation $M^2 \partial_t^2 \mathcal{E}(t,x) = 0$ for all $t > 0$. Consequently, the Poisson equation is satisfied for all times if and only if the deviation and its first derivative vanish at $t = 0$:
\begin{equation*}
    \mathcal{E}\vert_{t=0} = 0 \quad \text{and} \quad \partial_t \mathcal{E}\vert_{t=0} = 0,
\end{equation*}
which corresponds exactly to the initial conditions stated by equation \eqref{eq:init_cond}. 

Finally, it should be noted that the limit $\lambda \to 0$ inherently presupposes a strictly positive electron density ($n_e > 0$), which in turn guarantees the ellipticity of the reformulated equation in the quasi-neutral limit.
\end{proof}

\end{document}
